\documentclass[12pt,letterpaper]{article}

%Packages
\usepackage{amsmath}
\usepackage{amsthm}
\usepackage{amsfonts}
\usepackage{amssymb}
\usepackage{amscd}
\usepackage{dsfont}
\usepackage{enumerate}
\usepackage{fancyhdr}
\usepackage{mathrsfs}
\usepackage{bbm}
\usepackage{framed}
\usepackage{mdframed}
\usepackage{cancel}
\usepackage{float}
\usepackage{mathtools}
%\usepackage[]{mcode}
\usepackage{graphicx}
\usepackage{tikz}
\usepackage{hyperref}
\hypersetup{
    colorlinks=true,
    linkcolor=blue,
    filecolor=magenta,      
    urlcolor=cyan,
}
\usetikzlibrary{automata,arrows}

%Page formatting
\usepackage[letterpaper,voffset=-.5in,bmargin=3.5cm,footskip=1cm]{geometry}
\setlength{\parindent}{0.0in}
\setlength{\parskip}{0.1in}
\allowdisplaybreaks
\headheight 15pt
\headsep 10pt

% Common Commands
\newcommand\N{\mathbb N}
\newcommand\Z{\mathbb Z}
\newcommand\R{\mathbb R}
\newcommand\Q{\mathbb Q}
\newcommand\lcm{\operatorname{lcm}}
\newcommand\setbuilder[2]{\ensuremath{\left\{#1\;\middle|\;#2\right\}}}
\newcommand\E{\operatorname{E}}
\newcommand\V{\operatorname{V}}
\newcommand\Pow{\ensuremath{\operatorname{\mathcal{P}}}}

\DeclarePairedDelimiter\ceil{\lceil}{\rceil}
\DeclarePairedDelimiter\floor{\lfloor}{\rfloor}

\newcommand\definition{\textbf{Definition: }}
\newcommand\proposition{\textbf{Proposition: }}

% 22 style
\newcommand\hint[1]{\textbf{Hint}: #1}
\newcommand\note[1]{\textbf{Note}: #1}
\newenvironment{22enumerate}{\begin{enumerate}[a.]\itemsep0em}{\end{enumerate}}
\newenvironment{22itemize}{\begin{itemize}\itemsep0em}{\end{itemize}}
\fancypagestyle{firstpagestyle} {
  \renewcommand{\headrulewidth}{0pt}%
  \lhead{\textbf{CSCI 0220}}%
  \chead{\textbf{Discrete Structures and Probability}}%
  \rhead{Littman}%
}



\pagestyle{fancyplain}

\newcommand\EX{\mathds{E}}
\newcommand\PR{\mathds{P}}
\newcommand\zlam{Z_\lambda}
\DeclareMathOperator*{\argmin}{arg\,min}

%%%%%%%%%%%%%%%% COMPILE WITH \soltrue or \solfalse %%%%%%%%%%%%%%%%%%%%%%%%%%%%%
\newif\ifsol
\soltrue  % SOLUTIONS
% \solfalse % NO SOLUTIONS
%%%%%%%%%%%%%%%%%%%%%%%%%%%%%%%%%%%%%%%%%%%%%%%%%%%%%%%%%%%%%%

%%%%%%%%%%%%%%%% solution help %%%%%%%%%%%%%%%%%%%%%
\newcommand{\solu}[2]{ \begin{mdframed} \ifsol #2 \else \vspace{#1} \fi \end{mdframed} }
\newcommand{\solt}{ \ifsol (T) \fi }
\newcommand{\solf}{ \ifsol (F) \fi }
\newcommand{\solm}[1]{\ifsol \textit{(#1)} \fi}


\begin{document}
  \thispagestyle{firstpagestyle}
  \begin{center}
    {\large \textbf{Recitation 2}}
    
    {\large Circuits and Sets}

  \end{center}

	\section*{Part 1: Logic}
	
	In recitation 1, we started discussing logical equivalences and propositions. Recall some of the terminology for logic, most of which we talked about last week:
	
	\begin{22enumerate}
	
	\item A \textbf{propositional formula} is a function of one or more variables, each of which can be set to true or false, that evaluates to true or false. We call a propositional formula a \textit{proposition} for short. 
	
	\item The term \textbf{logical expression} is often used synonymously with the word proposition. 
	
	\item Two propositions are \textbf{logically equivalent} when they have  the same truth tables.
	
	\item A proposition is \textbf{valid} if it evaluates to true on any choice of inputs; it is true no matter what. That is, a valid proposition is logically equivalent to the expression ($p$ $\text{ OR }$ $\text{NOT } p$). This is also called a \textit{tautology}. 
	
	\item A proposition is \textbf{satisfiable} if it evaluates to true on some choice of inputs. A valid proposition is satisfiable, but so are many propositions which sometimes evaluate to false.
	
	\item If a proposition is not satisfiable, it evaluates to false on any choice of inputs; it is false no matter what. That is, it is logically equivalent to the expression ($p$ $\text{ AND }\text{NOT } p$). This is called a \textit{contradiction}.
	
	\end{22enumerate}

	\subsection*{From Propositions to Circuits}

Circuits are another way of representing truth tables. However, circuits can have multiple outputs, unlike propositions; that is circuits can represent more than one truth table. The number of outputs they have corresponds to the number of truth tables they represent.

Their representation is supposed to help visualize \textit{how} the output is computed given the inputs, and because computer science is very much about the \textit{how}, computer scientists often like these representations of truth tables very much. 

More specifically, the AND, OR, and NOT operations can be represented as follows:
\includegraphics[width=300px]{circuit.png}

Together, they can represent complicated logical expressions. For instance, the following circuit (created using Zoom whiteboard) is logically equivalent to the expression $\text{NOT } x_1 $ OR $ (x_2 $ AND $ x_3)$:

\includegraphics[width=200px]{circuit zoom example.png}

\textbf{Task 1}

Let's see if we can interpret these circuits! What are the logic expressions represented by these circuits?
\begin{enumerate}[a.]
    \item
    {\includegraphics[width=200px]{circuit_guess.png}}
    \begin{mdframed} \vspace{1cm} \end{mdframed}
    
    \item
    {\includegraphics[width=200px]{circuit_guess_2.png}}
    \begin{mdframed} \vspace{1cm} \end{mdframed}
\end{enumerate}

Now, let's practice drawing circuits with the following practice exercises. We recommend Google Jamboard, Google Slides, or sharing whiteboard on zoom to collaborate and draw together!

\textbf{Task 2}

\begin{enumerate}[a.]
        \item Draw a 3-input OR circuit using only 2-input OR gates
        \begin{mdframed} \vspace{3.5cm} \end{mdframed}
        
        \item \textit{Optional: }Draw a circuit for the XOR operator using AND, OR, and NOT gates.
        \begin{mdframed} \vspace{3.5cm} \end{mdframed}
\end{enumerate}

\subsection*{Checkpoint 1 — call over a TA!}

\newpage
\textbf{Task 3}

\begin{enumerate}[a.]

    \item Suppose you have a one-output circuit for some arbitrary truth table. Is there a quick way you could make a different circuit representing that same truth table?

    \textit{Hint:} Try adding gates to the end of the circuit. What gate can take in just one input?
    \begin{mdframed} \vspace{1cm} \end{mdframed}

    \item Using this method, are there infinitely many distinct one-output circuits representing that same truth table? How about infinitely many propositions?
    \begin{mdframed} \vspace{1cm} \end{mdframed}

    \item Suppose $C_1$ and $C_2$ are distinct one-output circuits but represent the same truth table. What reasons would we have to prefer one over the other? 

    \begin{mdframed} \vspace{1cm} \end{mdframed}

\end{enumerate}
\textbf{Task 4}

\begin{enumerate}[a.]
        \item Recreate/draw the half adder circuit from lecture, which calculates the addition of two single binary bits. More specifically, we have two inputs, A and B, as well as two outputs s (sum) and c (carry) that should produce the following truth table:
        
        \begin{center}
        {\includegraphics[width=150px]{half-adder2.png}}
        \end{center}
        
        For this part, be sure to use the XOR gate:
        \begin{center}
        {\includegraphics[width=100px]{XOR.png}}
        \end{center}
        
        \begin{mdframed} \vspace{3.5cm} \end{mdframed}
        
		\item \textit{Optional: }Given inputs $x_1$, $x_2$ and $x_3$, create a circuit which outputs true if and only if exactly one input is true.
		\begin{mdframed} \vspace{3.5cm} \end{mdframed}
			
		\item \textit{Optional: }Is the circuit you just created logically equivalent to $(x_1 \text{ XOR } x_2) \text{ XOR } x_3$?
		\begin{mdframed} \vspace{1cm} \end{mdframed}
		
\end{enumerate}

\subsection*{Checkpoint 2 - call over a TA!}

\newpage
      \section*{Set Theory}

      \textbf{Defn 1:} A \textbf{set} is a collection of objects with no repetition or order.

      \textbf{Defn 2:} $B$ is a \textbf{subset} of $A$ if every element in $B$ is also in $A$. This is written as $B \subseteq A$.

      \textbf{Defn 3:} The \textbf{integers} are the set $\Z = \{..., -2, -1, 0, 1, 2, ...\}$. The \textbf{non-negative integers} (also called the natural numbers) are the set $\N = \{0,1,2,...\}$.

      \textbf{Defn 4:} A number $n$ is \textbf{even} if $n = 2k$ for some $k \in \Z$. A number $n$ is \textbf{odd} if $n = 2k + 1$ for some $k \in \Z$.

      \textbf{Defn 5:} A number $n$ is \textbf{rational} if $n = \frac{a}{b}$ for some $a,b \in \Z$, where $b \neq 0$. 

	\textbf{Defn 6:} $\Pow{(A)}$, called the \textbf{power set} of $A$, is the set of all subsets of $A$.
	
	\textbf{Defn 7:} $A \cup B$ denotes the \textbf{union} of sets $A$ and $B$. This contains all the elements from $A$, and all of the elements from $B$. 
	
	\textbf{Defn 8:} $A \cap B$ denotes the \textbf{intersection} of sets $A$ and $B$. This contains only the elements that appear in both of the sets.
	
	\textbf{Defn 9:} $A - B$ denotes the \textbf{difference} of sets $A$ and $B$. This contains elements that appear in $A$ but not $B$.
	
	\textbf{Defn 10:} $\overline{A}$ denotes the \textbf{complement} of $A$ relative to some universal set $U$. $\bar{A} = U - A$, that is, it is everything except what is in $A$.
	
	\textbf{Defn 11:} $|A|$ denotes the \textbf{cardinality} of $A$, which is a count of the number of elements contained in $A$.
	
	\textbf{Defn 12:} The symbol $\forall$ means \textbf{for all}.
	
	\textbf{Defn 13:} The symbol $\exists$ means \textbf{there exists}.
	
	Answer the following questions. Discuss your answers!

    \textbf{Task 5}
      \begin{enumerate}[a.]
        \item True or False: $A$ is an arbitrary set. Answer true only if the statement is always true. That is, answer true only if for any possible set $A$, the statement is true. 
          \begin{enumerate}[i.]
            \item $A \subseteq A$
            \item $\{\} \subseteq A$
            \item $\{\} \in A$
		 \end{enumerate}
        \item True or False:  $\N \subseteq \Z$
        \item $\{0, 1, 9\} \subseteq \N$
        \item True or False: $\{-1.5, 9\} \subseteq \Z$
        
        \item Let $\mathbb{Q}$ be the set of rational numbers. 
        \begin{enumerate}[i.]
          \item $\mathbb{Q} \cap \N = \N$ 
          \item $\mathbb{Q} \cup \N = \R$
        \end{enumerate}
        
        \item True or False: $S$ is the set of planets in the solar system. $G$ is the set of planets in the galaxy. Jupiter is a student in the solar system.
        \begin{enumerate}[i.]
          \item $S \subseteq G$
          \item \text{Jupiter} $\subseteq S$
          \item \text{Jupiter} $\in S$
          \item $\{\text{Jupiter}\} \subseteq G$
        \end{enumerate}
        
	    \item True or False: If $A = \{1,2,4\}$ then $\{2,4\} \in \Pow{(A)}$
	    
        \item True or False: $A$ is a set, and $\Pow{(A)}$ is the set of all subsets of $A$. Answer true only if the statement is always true.
		\begin{enumerate}[i.]
		   \item $A \in \Pow{(A)}$
		   \item $A \subseteq \Pow{(A)}$
		   \item $\emptyset \in \Pow{(A)}$
		   \item $\emptyset \subseteq \Pow{(A)}$
          \end{enumerate}
          
        \item \textit{Optional:} In each of the following Venn diagrams, $A$, $B$, and $C$ are sets and are assumed to be subsets of a universal set (denoted by the rectangle). Write a set algebraic expression (i.e. one involving union, intersection, difference, and complement) in terms of $A$, $B$, and $C$ for each shaded in region. \begin{center}
          \includegraphics[scale = 0.5]{AuC.png} \quad \quad \quad \includegraphics[scale = 0.5]{notC.png}\\
          \includegraphics[scale = 0.5]{symdif.png}\\
          \end{center}
        
        \begin{mdframed} \vspace{2cm} \end{mdframed}
        
        
        \item \textit{Optional:} Call $a$ the cardinality of $A$ and $b$ the cardinality of $B$. Call $s$ the cardinality of $A $ $\cap$  $B$. For the third picture, what is the cardinality of the set formed from the expression you derived?
        
        \begin{mdframed} \vspace{1cm} \end{mdframed}
        
      \end{enumerate}

	\subsection*{Checkoff - Call over a TA!}

\end{document}
